\documentclass[10pt]{article}
\usepackage[margin=0.72in]{geometry}
\usepackage{amsmath,amssymb,mathtools,booktabs,longtable,array,microtype}
\usepackage[dvipsnames]{xcolor}
\usepackage[colorlinks=true,linkcolor=NavyBlue,urlcolor=BrickRed]{hyperref}
\usepackage{fancyhdr,enumitem}
\definecolor{ink}{HTML}{17212B}
\definecolor{navy}{HTML}{17365D}
\definecolor{teal}{HTML}{147D78}
\definecolor{pale}{HTML}{EEF4F7}
\color{ink}
\pagestyle{fancy}\fancyhf{}
\lhead{Proved matrix groups}\rhead{Formula compendium}\cfoot{\thepage}
\setlength{\headheight}{14pt}
\setlength{\parindent}{0pt}\setlength{\parskip}{4pt}
\setlist[itemize]{leftmargin=1.35em,itemsep=2pt,topsep=3pt}
\newcommand{\Sym}{\operatorname{Sym}}
\newcommand{\Tr}{\operatorname{Tr}}
\newcommand{\Fix}{\operatorname{Fix}}
\newcommand{\CT}{\operatorname{CT}}
\newcommand{\ExpS}{\operatorname{Exp}_{\sigma}}
\newcommand{\MGF}{\mathcal S}
\newcommand{\F}{\mathbb F}
\newcommand{\C}{\mathbb C}
\newcommand{\E}{\mathbb E}
\newcommand{\Aut}{\operatorname{Aut}}
\newcommand{\GL}{\operatorname{GL}}
\newcommand{\Cyc}{\operatorname{Cyc}}
\newcommand{\CTu}{\operatorname{CT}_u}
\newcommand{\resultbox}[1]{\begin{center}\fcolorbox{navy}{pale}{\parbox{0.93\linewidth}{#1}}\end{center}}
\title{\color{navy}\bfseries All New Matrix Groups and Formulas\\[3pt]
\large Consolidated theorem and verification compendium}
\author{Lambda Distributions workspace}
\date{20 July 2026}

\begin{document}
\maketitle

\begin{abstract}
This document is the cross-family index for the matrix-group results developed
in the workspace. It records every new group family, the formula actually
proved or corrected, its finite/stable/probabilistic status, and the principal
verification evidence. The detailed proofs, notebooks, and group-specific PDFs
remain in \path{proved_matrix_groups/packages} and
\path{proved_matrix_groups/documents}. The common object is the generalized
Molien or sigma-MGF series, but the compendium also includes two-sided unitary
series, nonuniform measures, heat kernels, and coefficientwise limits.
\end{abstract}

\resultbox{\textbf{Global outcome.} Exact finite-group formulas were checked by
explicit matrices, Reynolds ranks, orbit counts, or independent coefficient
recurrences. Compact-group formulas were checked by exact Weyl constant terms,
Lie-algebra kernels, or seeded Haar diagnostics. Claims requiring a correction
or a qualification are marked explicitly; no noncompact group is assigned a
normalized Haar probability.}

\tableofcontents

\section{Common notation and master identities}

For a representation $\rho:G\to GL(V)$ and a partition
$\tau=(\tau_1,\ldots,\tau_s)$, set
\[
 W_\tau(V)=\bigotimes_{a=1}^s\Sym^{\tau_a}V,
 \qquad |\tau|=\sum_a\tau_a.
\]
For finite $G$, the Reynolds projector gives the exact identity
\begin{equation}\label{eq:finite-master}
\boxed{\MGF_{G,V}(\mathbf t)=\frac1{|G|}\sum_{g\in G}
 \prod_i\det(I-t_i\rho(g))^{-1}
 =\sum_\tau\dim W_\tau(V)^G\,m_\tau.}
\end{equation}
Equivalently,
\[
 \MGF_{G,V}=\frac1{|G|}\sum_g
 \exp\!\left(\sum_{i,r\ge1}\frac{\Tr(\rho(g)^r)}r t_i^r\right).
\]
For compact $G$, normalized Haar integration replaces the finite average:
\begin{equation}\label{eq:compact-master}
\boxed{\MGF_{G,V}=\int_G\prod_i\det(I-t_i\rho(g))^{-1}dg
 =\sum_\lambda\dim(\mathbb S_\lambda V)^G s_\lambda
 =\sum_\tau\dim W_\tau(V)^Gm_\tau.}
\end{equation}
If $T$ is a maximal torus with representation weights $\mu$ and root system
$\Phi$, Weyl integration gives
\begin{equation}\label{eq:weyl}
\MGF_{G,V}=\frac1{|W|}\CT_z\left[
 \prod_{\alpha>0}(1-z^\alpha)(1-z^{-\alpha})
 \prod_{i,\mu}(1-t_i z^\mu)^{-m(\mu)}\right].
\end{equation}
For a probability measure $\nu$ that is not Haar, the coefficient is instead
the trace of its averaging operator:
\begin{equation}\label{eq:nonuniform}
[m_\tau]\MGF_{\nu,V}
=\Tr\!\left(\int_G\rho_{W_\tau}(g)d\nu(g)\right).
\end{equation}
Thus nonuniform coefficients need not be nonnegative integers.

\section{Finite abelian, monomial, dihedral, and wreath families}

\subsection{Finite abelian zero-weight theorem}
For $A=\prod_{\ell=1}^r C_{m_\ell}$ and
$V=\bigoplus_{j=1}^d\C_{w_j}$ with weight matrix $W=(w_{\ell j})$,
\begin{equation}\label{eq:abelian}
\boxed{\dim W_\tau(V)^A=\#\left\{q_{bj}\ge0:
 \sum_jq_{bj}=\tau_b,\quad
 \sum_{b,j}w_{\ell j}q_{bj}\equiv0\pmod{m_\ell}\ \forall\ell\right\}.}
\end{equation}
This is simultaneously a lattice count, a finite Fourier projection, and the
coefficient of the diagonal Molien average.

\subsection{Restricted monomial groups and $G(r,p,n)$}
For $G(r,p,n)$ on $\C^n$, $p\mid r$, the exact finite-$n$ series is
\begin{equation}\label{eq:grp}
\boxed{\MGF_{G(r,p,n)}=[u^n]\sum_{q=0}^{p-1}
 \ExpS\!\left(u\sum_{j\ge0}h_{qr/p+jr}\right).}
\end{equation}
This contains the Weyl groups $B_n/C_n=G(2,1,n)$ and $D_n=G(2,2,n)$.
More generally, for a code $A\le(\mathbb Z/r)^n$ stable under $H\le S_n$,
the coefficient for $D(A)\rtimes H$ is the average over $H$ of cycle-label
assignments whose phase-exponent vector lies in the dual code $A^\perp$.

\subsection{Wreath products}
For a finite matrix group $H\le GL(V)$ acting blockwise on $V^{\oplus n}$,
\begin{equation}\label{eq:wreath}
\boxed{\sum_{n\ge0}u^n\MGF_{H\wr S_n}(\mathbf t)
=\exp\!\left(\sum_{r\ge1}\frac{u^r}{r}
 \MGF_H(t_1^r,t_2^r,\ldots)\right).}
\end{equation}
The coefficientwise stable series is
\[
\MGF_{H\wr S_\infty}=\ExpS(\MGF_H-1),
\qquad \MGF_{H\wr S_n}\equiv\MGF_{H\wr S_\infty}\pmod{\deg>n}.
\]
The key block-cycle identity is
$\det(I-tg_C)=\det(I-t^{|C|}h_C)$.

\subsection{Generalized dihedral and dicyclic groups}
For $\operatorname{Dih}(A)=A\rtimes C_2$, the irreducibles are the two
extensions of every self-inverse character and the two-dimensional
$\operatorname{Ind}_A^G\chi$ for pairs $\{\chi,\chi^{-1}\}$. Substitution of
their diagonal/reflection spectra in \eqref{eq:finite-master} is the master
formula for arbitrary direct sums. All 33 generalized-dihedral checks passed.

For the dicyclic group $\operatorname{Dic}_m$ and its two-dimensional
irreducibles $\rho_k$, the series is the average of the $2m$ diagonal spectra
$(\zeta^{kj},\zeta^{-kj})$ and the $2m$ antidiagonal spectra. The coefficient
form is exactly $\dim W_\tau(\rho_k)^{\operatorname{Dic}_m}$; direct Reynolds,
character, and weight calculations agree. The regular-representation formula
for any finite group is
\begin{equation}\label{eq:regular}
\boxed{\MGF_{G,\mathrm{Reg}}=
 \sum_{d\ge1}\frac{N_d(G)}{|G|}\ExpS\!\left(\frac{|G|}{d}p_d\right),}
\end{equation}
where $N_d(G)$ counts elements of order $d$.

\section{Symmetric, alternating, cube, and permutation actions}

\subsection{Exact $S_n$ and $A_n$ formulas}
If $\lambda=1^{m_1}2^{m_2}\cdots\vdash n$, then the natural permutation
series is
\begin{equation}\label{eq:sn}
\boxed{\MGF_{S_n}=\sum_{\lambda\vdash n}\frac1{z_\lambda}
 \prod_{i,d}(1-t_i^d)^{-m_d(\lambda)}.}
\end{equation}
For $A_n$, retain the even cycle types and double their $S_n$ weights; the
extra orbit-splitting term appears only for all-distinct labels. In stable
range the $A_n$ and $S_n$ coefficients agree once two unused labels remain.

For the sign representation,
\[
\MGF_{S_n,\mathrm{sgn}}=\tfrac12\left[
\prod_i(1-t_i)^{-1}+\prod_i(1+t_i)^{-1}\right].
\]
For the standard module $V_n$, removing the trivial line gives
\[
\MGF_{S_n,V_n}=\sum_{\lambda\vdash n}\frac1{z_\lambda}
\prod_i\left[(1-t_i)\prod_{d\in\lambda}(1-t_i^d)^{-1}\right],
\qquad
\MGF_{S_n,V_n}\equiv\ExpS(h_2+h_3+\cdots)\pmod{\deg>n}.
\]

\subsection{Universal permutation-action formula}
For $G\curvearrowright X$ and vertex-cycle counts $c_d(g;X)$,
\begin{equation}\label{eq:perm}
\boxed{\MGF_{G,X}=\frac1{|G|}\sum_g\prod_{i,d}(1-t_i^d)^{-c_d(g;X)},
\quad [m_\tau]\MGF_{G,X}=\#\left(G\backslash\prod_j\Sym^{\tau_j}X\right).}
\end{equation}
Cycles may be reconstructed from fixed points by
\begin{equation}\label{eq:mobius}
\boxed{c_d(g;X)=\frac1d\sum_{r\mid d}\mu(d/r)\,|\Fix_X(g^r)|.}
\end{equation}
This single template covers Johnson $k$-subsets, Hamming schemes
$S_q\wr S_D$, Grassmann schemes $\operatorname{Gr}_q(N,k)$, polar and
dual-polar spaces, finite linear/projective/flag actions, and the projective
odd-symplectic actions below.

For $S_n$ on $k$-subsets, with $A_j(\sigma)$ the coordinate cycle counts,
\begin{equation}\label{eq:subsetfixed}
\boxed{|\Fix(\sigma^r)|=[u^k]\prod_{j=1}^{kr}
(1+u^{j/\gcd(j,r)})^{\gcd(j,r)A_j(\sigma)}.}
\end{equation}
For Hamming words, fixed words are the product of the numbers of alphabet
symbols fixed by the label product around each coordinate cycle. Grassmann
and polar actions replace this by the count of invariant subspaces, with the
isotropic/singular predicate in the polar case.

\subsection{The hyperoctahedral group on the binary cube}
For $B_n=\F_2^n\rtimes S_n$ acting affinely on $X_n=\F_2^n$, conditional on
$\pi\in S_n$ with $C(\pi)$ coordinate cycles,
\[
F_n(b,\pi)=|\Fix(b+\pi\cdot)|=
\begin{cases}2^{C(\pi)},&\text{with probability }2^{-C(\pi)},\\0,&\text{otherwise.}\end{cases}
\]
Therefore
\begin{equation}\label{eq:cubemoments}
\boxed{\mathbb E[F_n^m]=\frac{(2^{m-1})^{\overline n}}{n!}
=\binom{n+2^{m-1}-1}{n}.}
\end{equation}
Signed coordinate-cycle types $(a_\ell,b_\ell)$, their exact weights
$\prod_\ell((2\ell)^{a_\ell+b_\ell}a_\ell!b_\ell!)^{-1}$, the fixed-power
law, and \eqref{eq:mobius} give the full finite-$n$ cube series. In particular
$[m_{(2)}]=[m_{(1,1)}]=n+1$, so these coefficients do not stabilize.

\section{Finite affine, Lie-type, and exceptional finite families}

\subsection{Power characters, parabolic cycles, and bounded flags}
For every honest complex representation of a finite group, character values
on powers give the exact class-compressed formula
\begin{equation}\label{eq:power-character}
\boxed{\MGF_{G,V}=\sum_{[g]\subseteq G}\frac1{|C_G(g)|}
\exp\!\left(\sum_{r\ge1}\frac{\chi_V(g^r)}r p_r\right).}
\end{equation}
Thus the conjugacy classes, power maps, and character table determine the
entire Lambda-distribution.  In particular,
\begin{equation}\label{eq:power-coefficient}
\boxed{[m_\tau]\MGF_{G,V}=\frac1{|G|}\sum_{g\in G}\prod_j
\left(\sum_{\lambda\vdash\tau_j}\frac1{z_\lambda}
\prod_{r\ge1}\chi_V(g^r)^{m_r(\lambda)}\right).}
\end{equation}

For the coset action $X=G/H$, fixed points and cycles are
\begin{equation}\label{eq:parabolic-fixed-cycle}
\boxed{F_d(g;G/H)=\frac{|C_G(g^d)|}{|H|}|(g^d)^G\cap H|,\qquad
c_r(g;G/H)=\frac1r\sum_{d\mid r}\mu(r/d)F_d(g;G/H).}
\end{equation}
Substitution in \eqref{eq:perm} is an exact formula for every finite
Lie-type parabolic action $G/P$, including projective, Grassmann, polar, and
fixed-shape partial-flag actions.  For conjugation on $G$,
\begin{equation}\label{eq:finite-conjugation-cycles}
\boxed{c_r^{\rm conj}(g)=\frac1r\sum_{d\mid r}\mu(r/d)|C_G(g^d)|,\qquad
\MGF_{G,\rm conj}=\frac1{|G|}\sum_g\prod_{i,r}
(1-t_i^r)^{-c_r^{\rm conj}(g)}.}
\end{equation}

Fix $q$ and a flag shape $\mathbf d=(d_1<\cdots<d_s)$ with $d=d_s$.
For $X_n=\operatorname{Flag}_{\mathbf d}(\F_q^n)$ and $D=|\tau|$,
\begin{equation}\label{eq:bounded-flag-stability}
\boxed{[m_\tau]\MGF_{GL_n(q),X_n}\text{ is independent of $n$ for }n\ge dD.}
\end{equation}
The same range holds for $PGL_n(q)$; $SL_n(q)$ and $PSL_n(q)$ have the safe
range $n>dD$.  The stable coefficient is the sum, over $r\le dD$, of
$GL_r(q)$-orbits of spanning $\tau$-colored flag multisets.  Fixed-shape
isotropic and singular flags in each symplectic, unitary, or separate
orthogonal tower also stabilize after a finite tower-dependent threshold,
with the restricted form included in the supported configuration data.
Complete flags are excluded: $[m_{(1,1)}]=|W|$ already grows with rank.

Steinberg and specified genuine Weil representations are obtained by inserting
their characters in \eqref{eq:power-character}.  A Deligne--Lusztig virtual
character gives a completed Grothendieck-ring series whose coefficients can be
virtual.  A defining-characteristic linear action on $\mathfrak g(\F_q)$ is
not automatically a complex representation; its finite-set conjugation action
is an honest permutation alternative.

Exact verification used the standard/Steinberg module of
$S_3\cong GL_2(2)$, all $168$ matrices of $GL_3(2)$ on $G/P$, conjugation by
$S_3$, and consecutive stable ranks for projective points and $2$-planes.  All
26 coefficient and stability rows passed; the detailed proof and marimo
notebook are included in the proof anthology.

\subsection{Affine and finite classical actions}
For $\operatorname{AGL}_n(q)$ on $\F_q^n$,
\begin{equation}\label{eq:agl}
\boxed{\MGF_{n,q}=\frac1{q^n|GL_n(q)|}\sum_{A,b}
 \prod_r Q_r(\mathbf t)^{c_r(A,b)},\qquad
Q_r=\prod_i(1-t_i^r)^{-1}.}
\end{equation}
The fixed-point calculation is the linear system
$(I-A^r)x=(I+A+\cdots+A^{r-1})b$; coefficients stabilize for
$n\ge|\tau|-1$.

For nonzero vectors and projective points of finite linear groups,
\[
|\Fix_{\F_q^n\setminus0}(g^r)|=q^{\dim\ker(g^r-I)}-1,
\quad
|\Fix_{\mathbb P^{n-1}}(g^r)|=
\sum_{\lambda\in\F_q^\times}\frac{q^{\dim\ker(g^r-\lambda I)}-1}{q-1}.
\]
For complete flags, $[m_{(1,1)}]=|W|$, so rank stability already fails in
degree two. For stable polar point actions,
$[m_{(2)}]=[m_{(1,1)}]=3$ (coincident, perpendicular, nonperpendicular).

\subsection{$PSL_2(q)$, $\Delta$ groups, and code-monomial groups}
For $PSL_2(q)$ on $\mathbb P^1(\F_q)$, grouping identity, unipotent, split,
and nonsplit semisimple elements gives
\begin{align*}
\MGF=\frac1{|G|}\Big[&\ExpS((q+1)p_1)+(q^2-1)\ExpS(p_1+(q/p)p_p)\\
&+\frac{q(q+1)}2\sum_{\substack{m\mid(q-1)/d\\m>1}}\!\varphi(m)
\ExpS(2p_1+((q-1)/m)p_m)\\
&+\frac{q(q-1)}2\sum_{\substack{m\mid(q+1)/d\\m>1}}\!\varphi(m)
\ExpS(((q+1)/m)p_m)\Big],
\end{align*}
where $d=\gcd(2,q-1)$. The degree-three coefficient is $4+d$; the
cross-ratio term makes degree four grow with $q$.

For the standard $SU(3)$ monomial groups,
\begin{equation}\label{eq:delta}
\boxed{\MGF_{\Delta(3n^2)}=\tfrac13\mathcal A_n+\tfrac23\ExpS(p_3),
\quad
\MGF_{\Delta(6n^2)}=\tfrac16\mathcal A_n+\tfrac13\ExpS(p_3)+\tfrac12\mathcal B_n.}
\end{equation}
Here $\mathcal A_n$ averages the diagonal torus $N_n$, while $\mathcal B_n$
is the root-of-unity average of the transposition-coset factor
$((1-\eta^at_i^2)(1+\eta^{-a}t_i))^{-1}$.

The code-monomial family $D(C)\rtimes H$ obeys the dual-code cycle rule
described after \eqref{eq:grp}. Explicit groups of orders $8,24,18$ gave 57
exact comparisons. The $\Delta$ families gave 240 comparisons, and the
$PSL_2(q)$ actions gave 48.

\section{Braid, Pauli, extraspecial, Clifford, and Barnes--Wall groups}

\subsection{Finite braid images}
The verified images are scalar cyclic quotients of $B_3$, the permutation
quotient $B_3\twoheadrightarrow S_3$, and the finite Ising/Jones image. For a
scalar cyclic image $C_m$ on dimension $d$,
\[
[m_\tau]\MGF=\mathbf1_{m\mid|\tau|}\prod_i\binom{d+\tau_i-1}{\tau_i}.
\]
For the $S_3$ permutation image, the three class spectra give the exact
$1/6$ weighted identity/transposition/three-cycle formula. The Ising lift has
an eighth-root scalar selection rule; the projective adjoint removes scalar
dependence and is checked separately.

\subsection{Qudit Pauli and scalar-extended Pauli groups}
Let $q=p^f$ with odd $p$, $N=q^n$, and define
\[
A_\tau=\prod_i\binom{N+\tau_i-1}{\tau_i},\qquad
B_\tau=\mathbf1_{p\mid\tau_i\ \forall i}
\prod_i\binom{N/p+\tau_i/p-1}{\tau_i/p}.
\]
For the finite Pauli image,
\begin{equation}\label{eq:pauli}
\boxed{[m_\tau]\MGF_{P_{q,n}}=
\begin{cases}0,&p\nmid|\tau|,\\q^{-2n}A_\tau+(1-q^{-2n})B_\tau,&p\mid|\tau|.
\end{cases}}
\end{equation}
For $\widetilde P_{q,n}^{(s)}=\mu_sP_{q,n}$, replace the selection rule by
$s\mid|\tau|$. The noncentral term has a \emph{plus} sign; the proposed
qubit minus sign was disproved by the nonintegral coefficient it produced.

\subsection{Extraspecial and Clifford groups}
For an odd-$p$ extraspecial group of exponent $p$, the Pauli formula
\eqref{eq:pauli} applies. Exponent $p^2$ splits the noncentral sector according
to $g^p\in Z(E)$. For extraspecial $2$-groups of type
$\varepsilon\in\{+1,-1\}$, $N=2^n$,
\begin{align*}
\MGF_{E^\varepsilon}={}&2^{-2n}\Pi_2\prod_j(1-t_j)^{-N}\\
&+(\tfrac12+\varepsilon2^{-n-1}-2^{-2n})\prod_j(1-t_j^2)^{-N/2}
+(\tfrac12-\varepsilon2^{-n-1})\prod_j(1+t_j^2)^{-N/2}.
\end{align*}
This distinguishes $D_8$ from $Q_8$.

For a Clifford group containing all scalars, the one-sided series is exactly
$1$. The intrinsic object is the phase-neutral two-sided series
\[
[m_\alpha(\mathbf s)m_\beta(\mathbf t)]\mathcal B_G^0
=\dim(\Sym^\alpha V\otimes\Sym^\beta V^*)^G,
\qquad |\alpha|=|\beta|.
\]
Its $(1^k),(1^k)$ coefficient is the frame potential
$|G|^{-1}\sum_g|\Tr g|^{2k}$, so
\[
G\text{ is a unitary }k\text{-design}\iff
\Phi_k(G)=\sum_{\lambda\vdash k,\ \ell(\lambda)\le d}(f^\lambda)^2.
\]

\subsection{Real Clifford and Barnes--Wall layers}
For the real extraspecial Pauli group $E_m$ with $N=2^m$,
\begin{equation}\label{eq:realpauli}
\boxed{\begin{aligned}
\MGF_{E_m}=\frac1{2N^2}\Big[&\prod_j(1-t_j)^{-N}+\prod_j(1+t_j)^{-N}\\
&+(N^2+N-2)\prod_j(1-t_j^2)^{-N/2}\\
&+(N^2-N)\prod_j(1+t_j^2)^{-N/2}\Big].
\end{aligned}}
\end{equation}
For $\mathcal C_m=N_{O(2^m)}(E_m)$,
\[
\MGF_{\mathcal C_m}=\frac1{|O^+(2m,2)|}\sum_{A\in O^+(2m,2)}
 \frac1{|E_m|}\sum_{e\in E_m}\prod_j\det(I-t_je\widetilde A)^{-1}.
\]
The one-row stable coefficient $[m_{(2k)}]$ counts binary self-dual codes of
length $2k$ up to isomorphism. For the Barnes--Wall lattice group
$B_m\subset\mathcal C_m$ of index two,
\begin{equation}\label{eq:bw}
\boxed{\MGF_{B_m}=\MGF_{\mathcal C_m}+\MGF_{\mathcal C_m}^{\varepsilon_m}.}
\end{equation}
The twisted term is the lattice-sensitive sector.

\section{Compact classical, exceptional, and Schur-functor families}

\subsection{Defining representations of $SU(n)$, $SO(n)$, $O(n)$, and tori}
For the defining $SU(n)$ module,
\[
\dim W_\tau^{SU(n)}=\begin{cases}K_{(q^n),\tau},&|\tau|=nq,\\0,&n\nmid|\tau|.
\end{cases}
\]
For orthogonal groups,
\[
\dim W_\tau^{O(n)}=\sum_{\ell(2\delta)\le n}K_{2\delta,\tau},\qquad
\dim W_\tau^{SO(n)}=\dim W_\tau^{O(n)}+
\sum_{\ell(\delta)\le n}K_{2\delta+(1^n),\tau}.
\]
The row-length restriction is essential. For a weighted torus $T^r$,
\[
\MGF_{T^r,V}=\CT_{z_1,\ldots,z_r}\prod_{i,j}(1-t_i z^{w_j})^{-1}.
\]

\subsection{Exceptional and adjoint compact groups}
All use the exact Weyl constant term \eqref{eq:weyl}. Their verified
low-degree signatures are
\begin{align*}
\MGF_{G_2,7}&=1+h_2+e_3+O_{\ge4},\\
\MGF_{\operatorname{Spin}(7),8}&=1+h_2+h_4+s_{(2,2)}+e_4+O_{\ge6},\\
\MGF_{F_4,26}&=1+h_2+h_3+O_{\ge4},\\
\MGF_{E_6,27}&=1+h_3+O_{\ge6},\quad \MGF(t)=(1-t^3)^{-1},\\
\MGF_{E_7,56}&=1+e_2+h_4+s_{(2,2)}+e_4+O_{\ge6},\quad \MGF(t)=(1-t^4)^{-1},\\
\MGF_{E_8,248}&=1+h_2+e_3+O_{\ge4},\\
\MGF_{SU(2),\mathrm{Ad}}&=1+h_2+e_3+O_{\ge4},\\
\MGF_{SU(n),\mathrm{Ad}}&=1+h_2+h_3+e_3+O_{\ge4}\quad(n\ge3).
\end{align*}
The $E_6$ center imposes $3\mid|\tau|$; the centers of $E_7$ and
$\operatorname{Spin}(7)$ kill odd total degree. For $E_8$ the one-copy degrees
are $2,8,12,14,18,20,24,30$.

\subsection{Schur functors of $U(n)$, $O(n)$, and $USp(2n)$}
For every partition $\lambda\vdash d$,
\begin{equation}\label{eq:schurtrace}
\boxed{\Tr(S_\lambda(g)^r)=s_\lambda[X(g^r)]
=\sum_{\alpha\vdash d}\frac{\chi^\lambda(\alpha)}{z_\alpha}
\prod_{a\in\alpha}\Tr(g^{ra}).}
\end{equation}
For nonempty $\lambda$, the purely holomorphic $U(n)$ series is $1$; the
nontrivial series pairs the module with its conjugate. The stable orthogonal
and symplectic coefficients are the plethystic pushforwards
\[
[m_\tau]\MGF^O_{\lambda,\infty}=\langle\ExpS(h_2),h_\tau[s_\lambda]\rangle,
\qquad
[m_\tau]\MGF^{Sp}_{\lambda,\infty}=\langle\ExpS(e_2),h_\tau[s_\lambda]\rangle.
\]
In both cases $d|\tau|$ odd forces zero. The trace limits are polynomial
substitutions in Gaussian trace variables: circular complex for $U(n)$,
$\sqrt{k}A_k+\mathbf1_{2\mid k}$ for $O(n)$, and
$\sqrt{k}A_k-\mathbf1_{2\mid k}$ for $USp(2n)$.

\subsection{Local unitary tensor products and conjugation}
For $G=\prod_{a=1}^kU(d_a)$ on $V=\bigotimes_a\C^{d_a}$, the one-sided
series is $1$. The balanced coefficient vanishes unless
$|\alpha|=|\beta|$, and in stable bidegree
\begin{equation}\label{eq:local}
\boxed{\mathcal B_{G,V}\equiv\sum_\lambda z_\lambda^{k-2}
p_\lambda(\mathbf s)p_\lambda(\mathbf t).}
\end{equation}
Exactly,
$b_{(m),(m)}$ is the sum of squares of generalized Kronecker coefficients;
for two factors it counts partitions of $m$ obeying the dimension cutoffs.

For conjugation of $U(n)$ on $\operatorname{End}(\C^n)$,
\begin{equation}\label{eq:conj}
\boxed{\MGF=\int_{U(n)}\exp\!\left(\sum_{k\ge1}\frac{p_k}{k}
|\Tr(g^k)|^2\right)dg.}
\end{equation}
Finite coefficients are sums of squares of Littlewood--Richardson
multiplicities. For $n\ge|\tau|$ they count $H_\tau$-conjugacy orbits in
$S_{|\tau|}$, equivalently multisets of colored necklaces, giving the stable
Euler product $\prod_N(1-\operatorname{wt}N)^{-1}$.

\subsection{Proctor's odd symplectic group}
The complex stabilizer is
$W^*\rtimes(Sp_{2n}(\C)\times\C^\times)$ and is noncompact, so a normalized
Haar expectation does not exist. Its full algebraic covariant-invariant
series is $1$. The maximal compact replacements are
\[
\MGF_{Sp(n)\times U(1),\,W\oplus\chi}=\MGF_{Sp(n),W},\qquad
\MGF_{Sp(n),\,W\oplus\mathbf1}=\MGF_{Sp(n),W}\prod_i(1-t_i)^{-1}.
\]
Over $\F_q$, the honest complex permutation module on $\mathbb P(W\oplus L)$
obeys \eqref{eq:perm}; it has two point orbits, so $[m_{(1)}]=2$.

\section{Polyhedral, sporadic, and special 24-dimensional groups}

For a three-dimensional rotation through $\theta$, put
\[
\Phi_\theta=\prod_i\big((1-t_i)(1-2\cos\theta\,t_i+t_i^2)\big)^{-1}.
\]
The $A_4,S_4,A_5$ rotation series are the class-weighted averages of
$\Phi_\theta$ with their usual angle counts. Adjoining negatives gives the
$H_3$ reflection group; its one-variable Molien series is
$((1-t^2)(1-t^6)(1-t^{10}))^{-1}$. The binary lifts $2T,2O,2I\le SU(2)$
replace $\theta$ by paired half-angle spectra; $-I$ kills all odd degrees.
The tested $A_5$ modules include the permutation, deleted-permutation, and
three-dimensional rotation representations.

For the $M_{24}$ permutation action, if a class has cycle shape
$\prod_r r^{c_r(C)}$,
\[
\boxed{\MGF_{M_{24}}=\sum_C\frac{|C|}{|M_{24}|}
\prod_{i,r}(1-t_i^r)^{-c_r(C)}.}
\]
Five-transitivity gives agreement with $S_{24}$ through degree five; the first
extra orbit occurs at degree six. For the Leech representation of $Co_0$ with
frame shape $\prod_r r^{k_r(C)}$,
\[
\boxed{\MGF_{Co_0}=\sum_C\frac{|C|}{|Co_0|}
\prod_{i,r}(1-t_i^r)^{-k_r(C)},\qquad [m_\tau]=0\text{ for odd }|\tau|.}
\]
The signed-coordinate subgroup checks are not promoted to full-$Co_0$
coefficients.

\section{Nonuniform measures, heat flow, and limiting laws}

For a finite walk $\nu=\kappa^{*r}$, representation Fourier transforms obey
$\widehat\nu(W)=\widehat\kappa(W)^r$. For a cyclic diagonal representation,
\[
[m_\tau]\MGF_{\nu,V}=\sum_{q\bmod m}M_\tau(q)\widehat\nu(q).
\]
For a probability measure on $S_n$, the permutation series is the transformed
cycle-count PGF
$\mathbb E_\nu\prod_dQ_d(\mathbf t)^{C_d}$. This covers random transpositions,
random 3-cycles, adjacent transpositions, riffle shuffles, Ewens measures, and
general cycle weights. For Ewens parameter $\theta$,
\[
\MGF_{n,\theta}=\frac{[z^n]\exp(\theta\sum_{d\ge1}Q_dz^d/d)}
{[z^n]\exp(\theta\sum_{d\ge1}z^d/d)}.
\]

For $U(1)$ heat time $t$ and a weight decomposition
$W_\tau=\bigoplus_wM_\tau(w)\C_w$,
\begin{equation}\label{eq:heat}
\boxed{[m_\tau]\MGF_t=\sum_{w\in\mathbb Z}M_\tau(w)e^{-w^2t/2}.}
\end{equation}
This interpolates from dimension at $t=0$ to Haar invariants as
$t\to\infty$.

For fixed $k$, the $S_n$ action on $k$-subsets has the fixed-power formula
\eqref{eq:subsetfixed}. With independent $Z_j\sim\operatorname{Poisson}(1/j)$,
\[
F_{k,r}^{(\infty)}=[u^k]\prod_{j=1}^{kr}
(1+u^{j/\gcd(j,r)})^{\gcd(j,r)Z_j},
\quad
\MGF_\infty^{(k)}=\sum_\lambda\frac{\mathbb E\prod_aF_{k,\lambda_a}^{(\infty)}}
{z_\lambda}p_\lambda.
\]
For $k\ge2$ this limit is non-Poisson; its variance is $k$.

For a rooted tree whose branch isomorphism types $A$ occur with multiplicity
$m_A$,
\begin{equation}\label{eq:tree}
\boxed{\MGF_T(\mathbf t)=\prod_A[u_A^{m_A}]
\exp\!\left(\sum_{\ell\ge1}\frac{u_A^\ell}{\ell}\MGF_A(\mathbf t^\ell)\right).}
\end{equation}
This is exact for every finite tree. A general random-tree Lambda limit remains
a separate probabilistic problem.

\section{New-distribution formula and asymptotic ledger}

This section indexes every canonical note merged from \path{new_dists/}.
Each displayed expression is an explicit sigma-MGF (or the explicitly marked
projective/balanced replacement), rather than only a coefficient recipe.  The
full hypotheses, proofs, exceptional cases, and verification tables appear in
the final part of the proof anthology.

\subsection{Association-scheme permutation groups}
For any of the Johnson, Hamming, Doob, Grassmann, polar, building, translation,
or distance-regular actions $G\curvearrowright X$, define
\[
 F_r(g)=|\Fix_X(g^r)|,\qquad
 c_d(g)=\frac1d\sum_{r\mid d}\mu(d/r)F_r(g).
\]
Then the explicit group/action series is
\begin{equation}\label{eq:new-association}
 \boxed{\MGF_{G,X}=\frac1{|G|}\sum_{g\in G}
 \prod_{i,d\ge1}(1-t_i^d)^{-c_d(g)}.}
\end{equation}
For example, if $\lambda=1^{m_1}2^{m_2}\cdots\vdash v$,
\[
 F_r^{J(v,k)}(\lambda)=[u^k]\prod_{j\ge1}
 (1+u^{j/\gcd(j,r)})^{m_j\gcd(j,r)},
\quad
 \MGF_{J(v,k)}=\sum_{\lambda\vdash v}\frac1{z_\lambda}
 \prod_{i,d}(1-t_i^d)^{-c_d^J(\lambda)}.
\]
\paragraph{Asymptotics.}
For fixed $k$ and total degree $D$, Johnson and Grassmann coefficients stabilize
once the ambient rank is at least $kD$.  Hamming coefficients stabilize in the
alphabet size for $q\ge D$.  If both geometric parameters grow, pair-orbit
counts such as $\min(k,v-k)+1$, $d+1$, or $|W|$ can diverge, obstructing a raw
coefficientwise limit.

\subsection{Primitive-action families}
For an almost-simple coset action $G/M$, the substitution
\[
 F_r(g)=\frac{|C_G(g^r)|\,|(g^r)^G\cap M|}{|M|}
\]
in \eqref{eq:new-association} is the explicit MGF.  For diagonal, product,
twisted-wreath, holomorph-simple, and holomorph-compound actions, the same
series holds with the family-specific fixed-power products recorded in the
proof anthology.  In particular, for $H\wr S_n$ on $\Omega^n$,
\[
 F_r(h;\pi)=\prod_{C\in\Cyc(\pi)}
 |\Fix_\Omega(h_C^{r/e_C})|^{e_C},\qquad e_C=\gcd(|C|,r).
\]
\paragraph{Asymptotics.}
Almost-simple sequences require family data; product-action pair moments grow
polynomially, while diagonal-type moments grow by centralizer powers.  These
are generally multiset or critical-branching limits, not a universal ordinary
compound-Poisson law.

\subsection{Finite local-ring actions}
Let $R_a=\mathcal O/\pi^a\mathcal O$ with residue field $\F_q$, and let
$K_b(A)$ be the kernel size of $A$ over $R_b$.  The primitive-vector and
projective-point series for $\GL_n(R_a)$ are
\begin{align*}
 \MGF^{\rm prim}_{n,a}
 &=\frac1{|\GL_n(R_a)|}\sum_h\prod_{i,r}(1-t_i^r)^{-
 \frac1r\sum_{d\mid r}\mu(r/d)[K_a(h^d-I)-K_{a-1}(h^d-I)]},\\
 \MGF^{\rm proj}_{n,a}
 &=\frac1{|\GL_n(R_a)|}\sum_h\prod_{i,r}(1-t_i^r)^{-
 \frac1r\sum_{d\mid r}\mu(r/d)F_d^{\rm proj}(h)},\\
 F_d^{\rm proj}(h)
 &=\frac1{|R_a^\times|}\sum_{u\in R_a^\times}
 [K_a(h^d-uI)-K_{a-1}(h^d-uI)].
\end{align*}
\paragraph{Asymptotics.}
Fixed level and bounded degree give rank stability, but level growth does not:
$[m_{(1,1)}]\MGF^{\rm prim}=q^a$ and
$[m_{(1,1)}]\MGF^{\rm proj}=a+1$ in stable rank.

\subsection{Code-monomial matrix groups}
For $C\le(\mathbb Z/r)^N$, $P\le\Aut_{\rm perm}(C)$, and
$G(C,P)=D(C)\rtimes P$, write
\[
 H_a(\mathbf t)=\frac1r\sum_{s=0}^{r-1}\zeta^{-as}
 \prod_i(1-\zeta^st_i)^{-1}.
\]
The explicit weighted-cycle and dual-code forms are
\begin{equation}\label{eq:new-code}
 \boxed{\MGF_{C,P}=\frac1{|C||P|}\sum_{\pi,c}
 \prod_{i,Q\in\Cyc(\pi)}
 (1-\zeta^{\sum_{j\in Q}c_j}t_i^{|Q|})^{-1}
 =\frac1{|P|}\sum_\pi\sum_{\substack{u\in C^\perp\\u\text{ constant on }Q}}
 \prod_QH_{u_Q}(\mathbf t^{|Q|}).}
\end{equation}
For $P=1$, this is the complete weight enumerator
$\MGF_{D(C)}=\operatorname{cwe}_{C^\perp}(H_0,\ldots,H_{r-1})$.
\paragraph{Asymptotics.}
Length-growing families stabilize only under bounded-support orbit and
dual-code conditions.  Pure diagonal families already have
$[m_r]\MGF\ge N$, and many semidirect families have a diverging
$[m_{r,r}]$ coefficient.

\subsection{Multiplicative wreath and tensor-induced representations}
For $G_n=H\wr S_n$, an $H$-wreath type $\boldsymbol\mu$, and its power
transform $\boldsymbol\mu^{[r]}$,
\begin{equation}\label{eq:new-wreath-class}
 \boxed{\MGF_{G_n,W}=\sum_{\boldsymbol\mu\vdash_Hn}
 \frac1{Z_{\boldsymbol\mu}}
 \exp\!\left(\sum_{r\ge1}\frac{p_r}{r}
 \chi_W(\boldsymbol\mu^{[r]})\right).}
\end{equation}
For the product action on $\Omega^n$, use the fixed-power product in the
preceding subsection and \eqref{eq:new-association}.  For $V^{\otimes n}$,
\[
 \chi(g^r)=\prod_{Q\in\Cyc(\pi)}
 \chi_V(h_Q^{r/e_Q})^{e_Q},\qquad e_Q=\gcd(|Q|,r),
\]
inside \eqref{eq:new-wreath-class}.
\paragraph{Asymptotics.}
Additive block modules have compound-Poisson limits.  Product and tensor
modules usually do not: if $R_s=\#(H\backslash\Omega^s)$ then
$[m_{(1^s)}]=\binom{n+R_s-1}{R_s-1}$, and the analogous tensor coefficient
uses $a_s=\dim(V^{\otimes s})^H$.

\subsection{Central extensions and projective representations}
If $1\to Z\to\widetilde G\to G\to1$ and $Z$ acts through $\mu_r$, then
\begin{equation}\label{eq:new-central}
 \boxed{\MGF_{\widetilde G,V}=\frac1{|G|}\sum_{g\in G}\frac1r
 \sum_{\zeta\in\mu_r}\prod_i\det(I-\zeta t_i\rho(\widetilde g))^{-1},
 \quad [m_\tau]=0\ \text{unless }r\mid|\tau|.}
\end{equation}
For a genuinely projective representation, the lift-independent object is
\[
 \boxed{\mathcal B^{\rm proj}_{G,[V]}(\mathbf s,\mathbf t)=
 \CTu\int_G\prod_i\det(I-us_iU_g)^{-1}
 \prod_j\det(I-u^{-1}t_jU_g^{-*})^{-1}\,dg.}
\]
This covers Schur/spin covers, extraspecial central products, Weil lifts,
Clifford groups, determinant-one lifts, and sporadic covers.
\paragraph{Asymptotics.}
Finite scalar order gives congruence sectors; a continuous scalar center makes
the one-sided series equal to $1$.  Balanced design moments may stabilize,
while one-sided lift-dependent coefficients need not.

\subsection{Braid, TQFT, and quantum-image groups}
For a finite braid image $I$ with represented character $\chi$,
\[
 \boxed{\MGF_{I,V}=\sum_{[x]\subseteq I}\frac1{|C_I(x)|}
 \exp\!\left(\sum_{r\ge1}\frac{p_r}{r}\chi(x^r)\right).}
\]
For a monomial image $x=D(\alpha)P_\pi$, this becomes
$|I|^{-1}\sum_x\prod_{i,C}(1-\alpha_C(x)t_i^{|C|})^{-1}$.
Infinite images use the normalized Haar measure on the compact closure, not a
finite-image average.  The Ising lift has order $192$ and scalar center
$\mu_8$, so only total degrees divisible by eight survive.
\paragraph{Asymptotics.}
Dense $U(d)$ closure has one-sided series $1$; stable orthogonal and
symplectic closures are governed by contraction algebras.  Varying finite
images or fusion dimensions have no representation-free limit.

\subsection{Lattice automorphism groups}
For a positive-definite lattice $L$, write
$\det(I-tg)=\prod_m(1-t^m)^{a_m(g)}$.  Then
\begin{equation}\label{eq:new-lattice}
 \boxed{\MGF_L=\sum_{[g]\subseteq\Aut(L)}\frac1{|C_{\Aut(L)}(g)|}
 \prod_{i,m}(1-t_i^m)^{-a_m(g)}.}
\end{equation}
Since $-I\in\Aut(L)$, every odd-total-degree coefficient vanishes.
\paragraph{Asymptotics.}
Type $A$ has a Poisson-cycle limit; signed-permutation $B/C/D$ towers have a
stable signed-cycle law; repeated orthogonal sums have a marked
compound-Poisson limit.  Niemeier and isolated extremal/modular lattices are
finite collections, not intrinsic rank towers.

\subsection{Sporadic groups: the explicit $M_{11}$ series}
With $Z_\lambda=\prod_{i,r}(1-t_i^r)^{-c_r(\lambda)}$,
\[
 \boxed{\begin{aligned}
 \MGF_{M_{11},\C^{11}}=\frac1{7920}\big(&Z_{1^{11}}
 +165Z_{1^3\,2^4}+440Z_{1^2\,3^3}+990Z_{1^3\,4^2}\\
 &+1584Z_{1\,5^2}+1320Z_{2\,3\,6}
 +1980Z_{1\,2\,8}+1440Z_{11}\big).
 \end{aligned}}
\]
The deleted permutation module $W$ has the explicit series
\[
 \MGF_{M_{11},W}=\prod_i(1-t_i)\MGF_{M_{11},\C^{11}}.
\]
\paragraph{Asymptotics.}
There is no canonical sequence through the sporadic groups.  For a fixed
representation the formula is a rational function and ordinary coefficient
asymptotics apply; covers and other modules require their own represented
character and power maps.

\subsection{Solvable and finite-subgroup families}
For the natural monomial lamplighter representation
$L_{r,n}=\mu_r^n\rtimes C_n$, let
$A_{r,\ell}=r^{-1}\sum_{\zeta^r=1}\prod_i(1-\zeta t_i^\ell)^{-1}$.  Then
\[
 \boxed{\MGF^{\rm mon}_{L_{r,n}}=\frac1n\sum_{\ell\mid n}
 \varphi(\ell)A_{r,\ell}^{n/\ell}.}
\]
For $M(m,n,u)=\langle a,b\mid a^m=b^n=1,\ bab^{-1}=a^u\rangle$ in its
induced monomial representation,
\[
 \boxed{\MGF_{M(m,n,u)}=\frac1{mn}\sum_{x,y}
 \prod_{i,C\in\Cyc(y)}(1-\alpha_C(x,y)t_i^{|C|})^{-1}.}
\]
The affine lamplighter, Frobenius, unitriangular/Borel/parabolic, finite
unitary/symplectic, and reflection-wreath cases use
\eqref{eq:new-association}, the class-power formula, or
\eqref{eq:wreath} according to the stated representation.
\paragraph{Asymptotics.}
The monomial lamplighter has a linearly growing $[m_{r,r}]$ coefficient; the
affine action has necklace growth asymptotic to $r^n/n$.  Reflection wreath
blocks have a genuine compound-Poisson limit; isolated finite subgroups do not
form a canonical rank tower.

\subsection{Compact Spin and Pin groups}
For the type-$B_n$ spin module $\Delta_n$,
\begin{equation}\label{eq:new-spin}
 \boxed{\MGF_{\operatorname{Spin}(2n+1),\Delta_n}=
 \frac1{2^nn!}\CT_z\!\left[\mathcal D_{B_n}(z)
 \exp\!\left(\sum_{r\ge1}\frac{p_r}{r}
 \prod_{j=1}^n(z_j^{r/2}+z_j^{-r/2})\right)\right].}
\end{equation}
Type $D_n$ half-spin replaces the $B_n$ denominator and character by the
parity-selected half-spin weights.  Pin$^\pm(N)$ is the half-sum of the Spin
integral and the appropriate odd-component integral.
\paragraph{Asymptotics.}
$\dim(\Delta_n^{\otimes4})^{\operatorname{Spin}(2n+1)}=n+1$; half-spin fourth
moments also grow linearly by parity.  Hence the unnormalized spinor and
half-spin sigma-MGFs have no finite coefficientwise large-rank limit.

\subsection{Non-Haar Fourier-weighted laws}
For an arbitrary probability measure $\nu$ on a finite or compact group,
\begin{equation}\label{eq:new-nonhaar}
 \boxed{\MGF_{\nu,V}=\int_G\exp\!\left(\sum_{r\ge1}\frac{p_r}{r}
 \chi_V(g^r)\right)d\nu(g),\quad
 [m_\tau]\MGF_{\nu^{*k},V}=\sum_{\pi\in\widehat G}
 m_\pi(W_\tau)\Tr(\widehat\nu(\pi)^k).}
\end{equation}
For a central walk this reduces to
$\sum_\pi m_\pi(W_\tau)d_\pi\theta_\pi^k$.
\paragraph{Asymptotics.}
For fixed $G$, the largest nontrivial Fourier eigenvalue visible in $W_\tau$
gives the exact leading approach to Haar.  Heat kernels replace it by
$e^{-t\kappa_\pi/2}$; growing-rank cutoff claims require a fixed metric
normalization.

\subsection{Symmetric, spin-cover, Lie, Foulkes, and tree families}
For an $S_n$-module with Frobenius characteristic $f_n$,
\begin{equation}\label{eq:new-frobenius}
 \boxed{\MGF_{V_n}=\sum_{\mu\vdash n}\frac1{z_\mu}
 \exp\!\left(\sum_{r\ge1}\frac{p_r}{r}
 z_{\mu^{[r]}}[p_{\mu^{[r]}}]f_n\right).}
\end{equation}
This specializes to Specht, Schur-cover genuine, Foulkes, and free-Lie modules
(for the latter, $f_n=n^{-1}\sum_{d\mid n}\mu(d)p_d^{n/d}$).  Rooted-tree
leaf actions use the exact recursive cycle polynomial and the general
permutation MGF \eqref{eq:new-association}.
\paragraph{Asymptotics.}
Fixed-tail Specht modules have polynomial Poisson-chaos limits.  Growing
Foulkes shapes, free-Lie modules, and level-transitive rooted-tree actions have
diverging low moments and no raw coefficientwise limit; arbitrary random trees
need a separately specified probabilistic model.

\section{Verification ledger and corrections}

\begingroup\small
\setlength{\LTpre}{3pt}\setlength{\LTpost}{4pt}
\begin{longtable}{@{}p{0.25\linewidth}p{0.50\linewidth}r@{}}
\toprule
Family & Independent evidence & Reported checks\\\midrule
\endhead
Association schemes & matrices, cycle products, geometric fixed points, orbit traversal & 66\\
Matrix-group formula suite & compact kernels, finite Reynolds averages, spectral classes, dual-code counts & 928\\
Schur functors $U/O/USp$ & explicit Young-symmetrizer matrices versus plethystic traces & 120\\
$\Delta(3n^2),\Delta(6n^2)$ & explicit matrices versus $Q_n/R_n$ coefficient formulas & 240\\
Nonuniform $S_n$ & direct matrices versus cycle transform and Fourier operators & 267\\
$U(1)$ heat flow & quadrature versus weight/Casimir attenuation & 285\\
Hyperoctahedral cube & explicit $2^n$-dimensional matrices versus signed-cycle law & all $n\le5$\\
Rooted trees & recursive matrix closures versus wreath recurrence & 30\\
Generalized dihedral & Reynolds ranks versus classified spectral formula & 33\\
Odd symplectic & compact constant terms, vanishing certificate, finite permutation groups & 76\\
Other packages & package-specific exact/numerical suites documented in their PDFs & pass\\
\bottomrule
\end{longtable}
\endgroup

The consolidated work establishes the following corrections and boundaries:
\begin{itemize}
\item Orthogonal Schur shapes require $\ell(\lambda)\le n$.
\item A three-dimensional rotation factor is
$(1-t)^{-1}(1-2\cos\theta\,t+t^2)^{-1}$.
\item The scalar-extended qubit Pauli noncentral sector has a plus sign.
\item One-sided holomorphic series collapse to $1$ whenever a continuous
scalar center acts nontrivially; balanced two-sided series carry the content.
\item Proctor's complex odd symplectic group is noncompact and has no
normalized Haar probability.
\item Stable-range, finite-rank, numerical-Haar, and exact finite-group claims
are kept distinct throughout.
\end{itemize}

\paragraph{Artifact map.}
Sources live in \path{proved_matrix_groups/packages/},
\path{proved_matrix_groups/documents/}, and \path{new_dists/}.  All repository
marimo laboratories are centralized and cataloged under
\path{marimo_notebooks/}; the exact backends and pytest entry points remain
beside their group-family sources. The legacy paths
\path{for_this_guy/} and \path{output/pdf/for_this_guy/} are compatibility
links. This compendium indexes, but does not replace, the detailed proofs.

\end{document}
